\documentclass[11pt,a4paper]{article}
\usepackage[utf8]{inputenc}
\usepackage[russian,english]{babel}
\usepackage[T2A]{fontenc}
\usepackage{amsmath, amssymb, amsthm}
\usepackage{graphicx}
\usepackage{hyperref}
\usepackage{geometry}
\geometry{margin=2.5cm}

\title{Decoupled Observer Patch Holography (dOPH)\\ 
       \large Decoupled Observer Patch Holography}
\author{v04226079-sketch}
\date{April 27, 2026}

\begin{document}

\maketitle

\begin{abstract}
\textbf{Русская версия (для русскоязычных читателей)}

Представлена новая теоретическая модель \textbf{Decoupled Observer Patch Holography (dOPH)}. Ключевой параметр $P$ строго зафиксирован теоретическим значением $P \equiv \sqrt{8/3} \approx 1.63299$. Декуплинг наблюдателя реализован через строгий проекционный оператор. 

Модель демонстрирует высокую устойчивость на синтетических тестах (20\,000 галактик, стабильность 98.8\%) и успешно описывает реальную ротационную кривую Млечного Пути (reduced $\chi^2 = 1.84$). Подготовлен черновик анализа данных SPARC (175 галактик).

\vspace{1em}
\textbf{English Version}

We present a new theoretical model \textbf{Decoupled Observer Patch Holography (dOPH)}. The key parameter $P$ is strictly fixed to the theoretical value $P \equiv \sqrt{8/3} \approx 1.63299$. Observer decoupling is implemented via a strict projection operator.

The model shows high robustness in synthetic tests (20,000 galaxies, stability 98.8\%) and successfully reproduces the real Milky Way rotation curve (reduced $\chi^2 = 1.84$). A draft pipeline for the full SPARC dataset (175 galaxies) is prepared.
\end{abstract}

\section{Current Status of the Project / Текущий статус проекта}

\subsection{Русская версия (для русскоязычных читателей)}

К настоящему моменту в рамках проекта \textbf{dOPH} (Decoupled Observer Patch Holography) достигнуты следующие результаты:

\begin{itemize}
    \item Параметр $P$ переведён в строго теоретическое значение:
    \[
    P \equiv \sqrt{\frac{8}{3}} \approx 1.63299
    \]
    \item Декуплинг наблюдателя реализован строгим проекционным оператором.
    \item Проведены тяжёлые синтетические тесты на 20\,000 галактиках: средняя стабильность модели \textbf{98.8\%} (минимум 98.0\%), для LSB-галактик — \textbf{98.7\%}.
    \item Выполнен независимый тест на реальной ротационной кривой Млечного Пути (Eilers et al. 2019). Получено reduced $\chi^2 = \textbf{1.84}$ при фиксированном $P$.
    \item Подготовлен черновик `sparc_fittings_draft_v4.py` для фитинга полного набора данных SPARC (175 галактик). Файл \texttt{Rotmod\_LTG.zip} размещён в папке \texttt{data/sparc/}.
\end{itemize}

Модель демонстрирует высокую устойчивость на синтетических данных и приемлемые результаты на тесте Млечного Пути. Переход к реальным данным SPARC технически подготовлен, однако полный фитинг пока не запущен из-за ограничений текущего окружения (Android).

\subsection{English Version (for English-speaking readers)}

As of April 27, 2026, the following results have been achieved within the \textbf{dOPH} (Decoupled Observer Patch Holography) project:

\begin{itemize}
    \item The parameter $P$ has been fixed to its theoretical value:
    \[
    P \equiv \sqrt{\frac{8}{3}} \approx 1.63299
    \]
    \item Decoupling of the observer is implemented via a strict projection operator.
    \item Extensive synthetic tests were performed on 20,000 galaxies, yielding an average model stability of \textbf{98.8\%} (minimum 98.0\%), and \textbf{98.7\%} for low-surface-brightness (LSB) galaxies.
    \item An independent test was conducted on the real Milky Way rotation curve (Eilers et al. 2019). A reduced $\chi^2 = \textbf{1.84}$ was obtained with the fixed value of $P$.
    \item A draft fitting script `sparc_fittings_draft_v4.py` has been prepared for the full SPARC dataset (175 galaxies). The file \texttt{Rotmod\_LTG.zip} has been placed in the \texttt{data/sparc/} directory.
\end{itemize}

The model demonstrates high robustness on synthetic data and acceptable performance on the Milky Way test. 
The transition to real observational data from SPARC is technically prepared, although the full fitting has not yet been executed due to current computational environment limitations (Android). 
Further work will continue once a more suitable computing environment becomes available.

\section{Conclusion / Заключение}

Модель dOPH с фиксированным теоретическим параметром $P = \sqrt{8/3}$ показывает promising результаты на синтетических и реальных данных. Дальнейшая работа направлена на полный анализ SPARC, улучшение проекционного оператора и публикацию препринта.

The dOPH model with the theoretically fixed parameter $P = \sqrt{8/3}$ demonstrates promising results on both synthetic and real data. Future work will focus on the full SPARC analysis, refinement of the projection operator, and preparation of a preprint.

\end{document}
